\chapter{总结与展望}

\section{工作总结}

本文设计并实现了一套面向智能体的自进化强化学习系统，使智能体在无人工标注条件下自主产生训练数据、自主获得可信奖励、自主生成下一轮训练任务。主要工作如下：

（1）提出了跨步自进化机制。系统在每个训练步完成后，基于上一轮最优轨迹的执行概括与沙箱状态继承，自主生成下一轮训练任务。除初始种子外，后续所有训练查询由系统产生，形成不依赖外部数据的自进化闭环。

（2）设计了多智能体自主数据产生架构。执行者、观察者、提问者三个协同智能体使训练数据的产生、取证、追问全部自主化，无需人工标注。

（3）实现了超生-淘汰-选组的采样-训练解耦。推理超生、淘汰烂组、选梯度最优组训练，以"多推理少训练"提升每步训练信号质量，同时通过淘汰动态收缩超生量。

（4）构建了差分驱动的抗作弊奖励。奖励锚定在沙箱真实状态差分而非智能体自述，从取证机制上抑制奖励作弊。

（5）完成了训练基建与训练监控的工程实现。沙箱状态继承、推理训练异步分离、采样训练解耦集成保证自进化闭环物理运转；训练监控实时采集六项坍缩指标（奖励标准差、优势标准差、组内奖励标准差、策略 KL 散度、训练损失、存活组数），在五项终止判据触发时自动终止训练，保证自进化过程稳定可控。

系统在 Qwen3.5-9B 上端到端验证，完整自进化链路跑通（采样→奖励→淘汰→选组→GRPO 更新→训练监控→跨步种子生成），沙箱状态跨步继承验证通过，单元测试 239 项全部通过。结果表明，少量种子查询即可驱动完整强化学习训练，系统除初始种子外不依赖任何外部数据。

\section{研究展望}

本文以系统设计与工程搭建为主要贡献，未来可从以下方向深入：

（1）\textbf{自进化效率优化}。当前超生-淘汰-选组的淘汰阈值与选组策略为固定设计，未来可探索自适应调整策略，使系统根据训练监控指标动态优化采样与选组参数。

（2）\textbf{任务难度演化}。当前提问者生成的新种子未显式控制任务难度，未来可引入课程学习思想，使系统在自进化过程中逐步提升任务难度，实现更高效的策略进化。

（3）\textbf{更大规模验证}。本文在 16 卡上验证系统可行性，未来可在更大规模（64 卡及以上）上验证系统的可扩展性与长期自进化稳定性。

（4）\textbf{多模态扩展}。当前系统面向文本工具调用任务，未来可扩展至多模态智能体（含图像、视频），使自进化机制适用于更广泛的智能体场景。
